% =================
% Copyright + Aviso ao Leitor — legível (raggedright), uma página
% =================
\cleardoublepage
\thispagestyle{empty}
{\small
\setlength{\parindent}{0em}
\setlength{\parskip}{0.7em}
~
\vspace{0.3em}
\vspace{1em}

Primeira edição, \editionyear{}\\
\publishercity{}: \publisher{}

Copyright \copyright{} \editionyear{} \authorname{} e \coauthorname{}

Todos os direitos reservados. Nenhuma parte deste livro pode ser reproduzida
por qualquer meio sem a autorização prévia e por escrito do autor, exceto
para citações breves em resenhas e artigos.

\textbf{ISBN:} \isbn{}. \quad Impresso no Brasil.

\textbf{Ficha catalográfica (CIP):} ver página seguinte.

\vspace{0.4em}
{\footnotesize
\textbf{Aviso ao leitor.} Este romance é uma obra de ficção baseada em
eventos históricos documentados. Os campos de concentração do Ceará (1915),
o Hospital Colônia de Barbacena e as cerca de 60.000 mortes em suas
dependências são fatos reais. Joaquim Antônio Correia, o Dr. Heitor Oliveira
e a Dra. Lúcia Mendes são personagens ficcionais. O ciclo de contaminação
descrito neste livro é uma construção narrativa.

\textit{Qualquer semelhança com a sua experiência de leitura é inteiramente
intencional.}

O autor se isenta de qualquer responsabilidade por noites mal dormidas,
cheiros adocicados inexplicáveis ou olhos amarelos no canto do quarto.}

\vfill
\begin{center}
\includegraphics[width=0.16\linewidth]{frontmatter/logo-black.png}\\[0.4em]
{\scshape Publicado por \publisher{} \quad·\quad \editionyear{}}
\end{center}
}
\clearpage
